\documentclass[11pt]{article}
\usepackage[margin=1in]{geometry}
\usepackage{amsmath,amssymb}
\usepackage[hidelinks]{hyperref}

\title{BCH codes of length seven have different covering radii\\[2pt]
  \large Disproof of Conjecture 00000001017}
\author{}
\date{}

\begin{document}
\maketitle

\noindent
Under the standard primitive narrow-sense binary BCH convention,
the designed-distance-three code of length $2^3-1=7$ has covering radius
one. The proposed value is $2^{\lceil 3/2\rceil}-1=3$. Another code under
the same BCH convention and at the same length has radius three, so the
length alone does not determine the radius.

\section*{Designed distance three}

Let $m\ge2$, let $n=2^m-1$, and choose a generator $\alpha$ of the
multiplicative group of $\mathbb F_{2^m}$. Identify a binary word
$(w_0,\ldots,w_{n-1})$ with $w(x)=\sum_{i=0}^{n-1}w_i x^i$. The primitive
narrow-sense BCH code of designed distance $\delta$ consists of the words
that vanish at $\alpha,\alpha^2,\ldots,\alpha^{\delta-1}$. For a binary
code $C$, its covering radius is
\[
\rho(C)=\max_{w\in\mathbb F_2^n}\min_{c\in C}d_H(w,c),
\]
where $d_H$ is Hamming distance.

At designed distance three, $w(\alpha^2)=w(\alpha)^2$ because the
coefficients are binary. Consequently the word belongs to the code exactly
when its syndrome $w(\alpha)$ is zero. If the syndrome is nonzero, it is
$\alpha^i$ for some $0\le i<n$. Flipping coordinate $i$ replaces $w(x)$ by
$w(x)+x^i$, which has syndrome zero. This gives covering radius at most one.
The word represented by $w(x)=1$ has nonzero syndrome, so the radius is
exactly one. In particular, the proposed expression is at least three for
every $m\ge3$ and fails on this family throughout that range.

At $m=3$ the field can be written as
$\mathbb F_2[x]/(x^3+x+1)$ with $\alpha$ equal to the residue class of $x$.
The cubic $g(x)=x^3+x+1$ has no root in $\mathbb F_2$ and is irreducible.
The element $\alpha$ is nonzero and differs from $1$, so it has order
seven. The condition $w(\alpha)=0$ is equivalent to divisibility by
$g(x)$. Thus the code is generated by $g(x)$ and consists of the products
$q(x)g(x)$ with $\deg q\le3$, including the zero polynomial. Distinct
nonzero powers of $\alpha$ rule out zero-syndrome words of weight one or
two, and $g(x)$ has weight three. This is the cyclic binary $[7,4,3]$
Hamming code~\cite{tervo}.

\section*{A second code of length seven}

At designed distance seven the same BCH construction requires all six
nonidentity powers of $\alpha$ as roots. The monic polynomial with those
roots is
\[
\frac{x^7-1}{x-1}=1+x+x^2+x^3+x^4+x^5+x^6.
\]
It follows that the degree-at-most-six binary words satisfying the root
conditions are the zero word and the all-ones word. Designed distance
five already gives the same code under this convention. It includes
$\alpha^3$ as a root, so the Frobenius map gives the roots $\alpha^6$ and
$\alpha^{12}=\alpha^5$ as well.

The distance from a word of weight $s$ to this repetition code is
$\min(s,7-s)$. The maximum is three, attained at weights three and four.
The radius therefore differs from that of the designed-distance-three
code, although both block lengths are seven.

\section*{Lean verification}

The Lean 4 project represents seven-bit words as binary polynomials and
multiplies by the two explicit generators. It evaluates the required
roots with the relation $\alpha^3=\alpha+1$ and checks that the root-zero
words coincide with the corresponding encoder images. It checks the
upper covering bounds for all seven-bit words and checks witnesses that
rule out the next smaller radius in each code. The theorem
\texttt{C1017.bch\_counterexample} states both exact radii and the failure of
the proposed value for designed distance three directly in terms of the
BCH root conditions. Run \texttt{lake build} in the \texttt{lean}
directory to verify it.

\begin{thebibliography}{1}
\bibitem{tervo}
R.~Tervo, \emph{BCH Code Generator}, ECE 4253 Digital Communications,
University of New Brunswick.
\href{https://www.ece.unb.ca/cgi-bin/tervo/bch.pl}{Explicit length-seven example}.
\end{thebibliography}

\end{document}
